The introduction has already named the foundational threads on which \cmapagent\ rests; this section situates those threads within the broader landscape and identifies the converging research directions on which the system's design draws.

\subsection{Retrieval augmentation, tool use, and grounded LLM agents}
Two complementary mechanisms have emerged for grounding large language models in external resources.
Retrieval-augmented generation (RAG) couples parametric language models with an explicit, updateable external memory, improving factuality and enabling provenance in knowledge-intensive settings \citep{Lewis2020RAG}; in scientific contexts, agentic variants iterate retrieval and synthesis across full-text papers and have been shown to match or exceed expert performance on literature-research tasks \citep{Lala2023PaperQA,Skarlinski2024PaperQA2}.
Tool augmentation extends this grounding from memory to action: ReAct interleaves reasoning and acting, demonstrating benefits for interpretability and reduced hallucination when models can query external resources \citep{Yao2022ReAct}; Toolformer shows that models can learn when and how to invoke tools through self-supervision \citep{Schick2023Toolformer}; and modular neuro-symbolic architectures such as MRKL argue for combining large language components with specialized reasoning and knowledge modules \citep{Karpas2022MRKL}.
\cmapagent\ draws on both threads: the knowledge base provides updateable retrieval-based grounding, while an explicit tool layer turns retrieved context into validated, executable operations.

\subsection{Scientific domain agents and their evaluation}
Recent systems demonstrate the scientific potential of tool-augmented LLM agents across domains.
ChemCrow integrates expert tools for chemistry and reports multi-step task success across synthesis and discovery workflows \citep{Bran2024ChemCrow}, and Coscientist demonstrates an agentic system that combines planning, tool use, and execution to carry out multi-step experimental procedures in an automated laboratory setting \citep{Boiko2023Coscientist}.
In ocean science, domain-focused models such as OceanGPT aim to improve performance on oceanography tasks via instruction data and domain benchmarks \citep{Bi2023OceanGPT}; in the geospatial domain, GeoGPT exemplifies an autonomous workflow paradigm in which a language model plans and executes GIS tools sequentially to complete user-specified geospatial tasks \citep{Zhang2024GeoGPT}; and beyond single-agent systems, multi-agent frameworks have been proposed for hypothesis generation and scientific exploration, including graph- and knowledge-based approaches \citep{Ghafarollahi2024SciAgents}.
A complementary strategy improves domain knowledge through pretraining: domain-specific foundation models for geoscience and climate (e.g., K2 and ClimateGPT) often pair this with retrieval augmentation for factuality \citep{Deng2023K2,Thulke2024ClimateGPT}.
Recent reviews of AI for geoscience synthesize progress in large-scale models, physics-informed approaches, and geoscientific applications, while emphasizing data governance, uncertainty, and interpretability challenges \citep{Zhao2024AIforGeoscience}; agentic RAG systems for environmental data can be viewed as one response to these challenges, grounding domain language in curated metadata, routing computations through validated tools, and returning artifacts that support auditability.
Reliable evaluation of such systems remains difficult because they must plan, act, recover from errors, and manage multi-step interactions: general-purpose agent benchmarks have been proposed for interactive environments \citep{Liu2023AgentBench}; tool-oriented benchmarks such as ToolBench (introduced alongside ToolLLM) and StableToolBench focus on large-scale evaluation of tool use and the stability of tool APIs \citep{Qin2023ToolLLM,Guo2024StableToolBench}; and for data analysis agents specifically, DSEval and IDA-Bench emphasize the end-to-end lifecycle and multi-round nature of practical analysis workflows \citep{Zhang2024DSEval,Li2025IDABench}.
These lines of work motivate \cmapagent\ as a domain-grounded, tool-executing assistant focused on scientific data discovery and analysis rather than model pretraining alone.

\subsection{Natural-language interfaces to scientific data}
Natural-language interfaces to structured data have a long history, and LLMs have renewed interest in text-to-SQL and database interfaces.
However, scientific workflows typically require more than query translation; they demand validation, iterative refinement, visualization, and the production of reusable analysis artifacts.
\cmapagent\ is positioned as an agentic layer that combines grounded catalog retrieval with domain tools and artifact-oriented outputs for ocean and environmental data analysis.
